\documentclass{article}
\title{Differential Equations - Notes}
\author{Arun Felix}
\begin{document}
\maketitle
\tableofcontents
\newpage
\part{Basic Concepts}
\section{Definitions}

\subsection{What is a differential equation?}

\begin{itemize}
    \item A differential equation is an equation that contains derivatives, either ordinary or partial derivatives
    \item A popular example is newtons second law of motion, $F = ma$
    \item This is a differential equation since $a = \frac{dv}{dt}$
    \item Given this, we can actually write Newton's second law as a differential equation: $m \frac{dv}{dt} = F(t,v)$
\end{itemize}

\subsection{What is Order?}

\begin{itemize}
    \item The order is the largest derivative present in the differential equation.
    \item For example, $\frac{\partial ^3 u}{\partial x^2 \partial t} = 1 + \frac{\partial u}{\partial y}$ is a third order differential equation since it's greatest derivative is 3.
\end{itemize}

\subsection{Ordinary vs. Partial Differential Equation}

\begin{itemize}
    \item An ordinary differential equation only has original differential equations, ode.
    \item A partial differential equation will have partial derivatives in it.
\end{itemize}

\subsection{Linear Differential Equation}

A linear differential equation is a differential equation following the form:

\begin{center}
    $a_n (t) y^{(n)}(t) + a_{n - 1} (t) y^{(n-1)} + ... + a_1 (t) y' (t) + a_0 (t) y (t) = g(t)$ 
\end{center}

\subsection{Initial Conditions}

An initial condition allows us to determine which solution we're after. For example, $y(t_0) = y_0$ or $y^{(k)}(t_0) = y_k$.

\subsection{Implicit vs. Explicit Solutions}

An implicit solution is a solution given in the form of $y = y(t)$ \\
An explicit solution is a solution that isn't implicit. 

\part{First Order Differential Equations}
\section{Linear Differential Equations}
\subsection{Introduction}
A linear first order differential equation is an equation in the form:

\begin{center}
    $\frac{dy}{dt} + p(t)y = g(t)$
\end{center}

\begin{itemize}
    \item Here, we can assume $p(t)$ and $g(t)$ are continuous.
\end{itemize}

Now, let's assume the existence of a variable, $\mu (t)$, where we multiple all factors by it

\begin{center}
    $\mu(t)\frac{dy}{dt} + \mu(t)p(t)y = \mu(t)g(t)$
\end{center}

To solve problems using this:
\begin{itemize}
    \item Put the differential equation in the correct form above.
    \item Find the integrating factor, $\mu (t)$, where $\mu (t) = e^{\int p(t) dt}$
    \item Multiply everything in the differential equation with $\mu (t)$, and verify the leftside becomes the product rule, $(\mu(t) y(t))'$.
    \item Integrate both sides.
    \item Solve for $y(t)$.
\end{itemize}

\subsection{Example 1}

Let's try to find the solution to the following differential equation:

\begin{center}
    $\frac{dv}{dt} = 9.8 - 0.196v$
\end{center}

We can see that we need to put the equation into linear form. 

\begin{center}
    $\frac{dv}{dt} + 0.196v = 9.8$
\end{center}

Let's solve for $\mu (t)$

\begin{center}
    $\mu (t) = e^{\int 0.196 dt} = e^{0.196t}$
\end{center}

Next, we can plug in $\mu (t)$ in all variables

\begin{center}
    $e^{0.196t}\frac{dv}{dt} + e^{0.196t} 0.196v = 9.8e^{0.196t}$
\end{center}

We can rearrange some numbers:

\begin{center}
    $e^{0.196t}\frac{dv}{dt} +  0.196e^{0.196t}v = 9.8e^{0.196t}$
\end{center}

this resembles the product rule, hence we can say:

\begin{center}
    $(e^{0.196t} v)' = 9.8e^{0.196t}$
\end{center}

Now, we can integrate:

\begin{center}
    $\int(e^{0.196t} v)' dt = \int9.8e^{0.196t} dt$
\end{center}

\begin{center}
    $e^{0.196t} v = 9.8\int e^{0.196t}dt$ 
\end{center}

for the right equation, we can use u sub

\begin{center}
    $e^{u} dt$, where $u = 0.196t$, $du = 0.196dt$, and $dt = \frac{du}{0.196}$
\end{center}

We can plug this into our equation

\begin{center}
    $e^{0.196t} v + k = 9.8 \cdot \frac{1}{0.196} \int e^{u} du$ \\
    $e^{0.196t} v + k = 50 \int e^{u} du$ \\
    $e^{0.196t} v + k = 50 e^{0.196t} + c$
\end{center}

Finally, we can solve for y, or v in this context:

\begin{center}
    $ v = \frac{50 e^{0.196t} + c}{e^{0.196t}}$
\end{center}

\section{Separable Equations}
\subsection{Introduction}
A seperable differential equation is an equation where you can separate the variables:
\begin{center}
    $N(y) \frac{dy}{dx} = M(x)$
\end{center}

To put it simply, all the ys need to be on one side, xs on the other. To solve such a differential equation, you'd integrate both sides.

\begin{center}
    $\int N(y) \frac{dy}{dx} dx = \int M(x) dx$
\end{center}

We can use substitution to get this instead:

\begin{center}
    $\int N(u) du = \int M(x) dx$
\end{center}

\subsection{Example}

Solve the following differential equation and determine the interval of validity for the solution.

\begin{center}
$\frac{dy}{dx} = 6y^2 x$, and $y(1) = \frac{1}{25}$
\end{center}

The first step is seperating the Xs and Ys

\begin{center}
    $dx \cdot \frac{dy}{dx} = 6y^2 x \cdot dx$ \\ 
    $\frac{dy}{y^2} = \frac{6y^2 x \cdot dx}{y^2}$ \\
    $\frac{1}{y^2} dy = 6x dx$
\end{center}

From here, we can do integration:

\begin{center}
    $\int \frac{1}{y^2} dy = \int 6x dx$ \\
    $\int y^{-2} dy = 3x^2 + C$ \\
    $-y^{-1} = 3x^2 + C$ \\
    $-\frac{1}{y} = 3x^2 + C$
\end{center}

now, we can solve for y:

\begin{center}
    $y = \frac{1}{3x^2 + C}$
\end{center}

Finally, we can solve for C to get our differential equation. This is trivial, as we have $y(1) = \frac{1}{25}$

\begin{center}
    $\frac{1}{25} = \frac{1}{3 + C}$
\end{center}

We can easily figure out that $C = 22$.

Hence our final differential equation is:

\begin{center}
    $y = \frac{1}{3x^2 + 22}$
\end{center}


\section{Exact Equations}
\subsection{Introduction}
An exact equation follows the form: $M(x,y) + N(x,y) \frac{dy}{dx} = 0$

Let's create two variables:

\begin{itemize}
    \item $\psi_x = M(x,y)$
    \item $\psi_y = N(x,y)$
\end{itemize}

Plugging it into our equation, we get:

\begin{center}
$\psi_x + \phi_y \frac{dy}{dx} = 0$
\end{center}

We can then use the chain rule to find the derivative of the differential equation:

\begin{center}
    $\frac{d}{dx} (\psi (x,y(x))) = 0$
\end{center}
this gives us an implicit solution: $\psi(x,y) = c$.

\subsection{What is a exact equation}

An exact equation is an equation, given $M_y$ and $N_x$, if their partial derivatives are equal to eachother.

\subsection{Example}

Solve the following IVP and find the interval of validity for the solution

$2xy - 9x^2 + (2y + x^2 + 1)\frac{dy}{dx} = 0$, where $y(0) = -3$.

\begin{itemize}
    \item We can set two variables.
    \item $M(x,y) = 2xy - 9x^2$
    \item $N(x,y) = 2y + x^2 + 1$
\end{itemize}

The first thing we can do is check if this is an exact equation:

\begin{center}
    $\frac{\partial M}{\partial y} = 2x$ \\
    $\frac{\partial N}{\partial x} = 2x$
\end{center}

We can see the partial derivatives are identical, hence it's an exact equation.

Next, we need to integrate one of the equations:

\begin{center}
    $\int M(x,y) = \int 2xy - 9x^2 dx = x^2 y - 3x^3 + h(y) $
\end{center}

Notice how, instead of a constant, we use $h(y)$.

We actually have $h'(y)$, it's $N = 2y + x^2 + 1$.

Now, we find $\frac{\partial M}{\partial y}$ the integral of M, giving us $x^2 + h'(y)$

We can set $\frac{\partial M}{\partial y} = M$, or $x^2 + h'(y) = 2y + x^2 + 1$.

\begin{center}
    $x^2 + h'(y) - x^2 = 2y + x^2 + 1 - x^2$ \\
    $h'(y) = 2y + 1$
\end{center}

We can integrate $h'(y)$ to get $h(y)$.

\begin{center}
    $\int 2y + 1 dy = y^2 + y + k$
\end{center}

We can plug this into our equation:

\begin{center}
    $\psi (x,y) = x^2 y - 3x^3 + y^2 + y + k$
\end{center}

Now, we need to solve for k, or c

\begin{center}
    $\psi (x,y) = x^2 y - 3x^3 + y^2 + y + k = c$ \\
    $\psi (x,y) = x^2 y - 3x^3 + y^2 + y + k - k = c - k$ \\
    $\psi (x,y) = x^2 y - 3x^3 + y^2 + y = c$ \\
    $\psi (x,y) = x^2 y - 3x^3 + y^2 + y = c$ Here, we use $y(0) = -3$ \\
    $\psi (x,y) = (0)^2 (-3) - 3(0)^3 + (-3)^2 + (-3) = c$ Here, we use $y(0) = -3$ \\
    $\psi (x,y) = 9 -3 = c$, where $c = 6$ \\
    hence our solution is $\psi (x,y) = x^2 y - 3x^3 + y^2 + y - 6 = 0$ \\
\end{center}

\section{Bernoulli Differential Equation}
\subsection{Introduction}
Bernoulli Equations are equations that follow the form:

\begin{center}
    $y' + p(x)y = q(x)y^n$
\end{center}

Here, we assume $p(x)$ and $q(x)$ are continuous functions, and n is a number.

if n = 0 or 1, we'll have a linear equation, which is easy to solve.

if n isn't 0 or 1, we need to divide both sides by $y^n$, substitute $v = y^{1 - n}$. Once we do this, we'll get a linear equation we can solve.

\subsection{Example}

Solve the following IVP and find the interval of validity for the solution:

\begin{center}
    $y' + \frac{4}{x} y = x^3 y^2$, where $y(2) = -1$ and $x > 0$
\end{center}

We can see our differential equation resembles the bernoulli equation, where:

\begin{itemize}
    \item $p(x) = \frac{4}{x}$
    \item $q(x) = x^3$
    \item $n = 2$
\end{itemize}

now, we need to divide both sides by $y^n$

\begin{center}
    $\frac{y' + \frac{4}{x} y}{y^2} = \frac{x^3 y^2}{y^2}$ \\
    $\frac{y' + \frac{4}{x} y}{y^2} = \frac{x^3 y^2}{y^2}$ \\
    $y^{-2} y' + \frac{4}{x} y^{-1} = x^3$
\end{center}


From here, we can substitute, using $v = y^{1 - n} = y^{1 - 2} = y^{-1}$ and $v' = -y^{-2} y'$, giving us:

\begin{center}
    $-v' + \frac{4}{x} v = x^3$
\end{center}

Because $\frac{dv}{dt}$ needs to be positive, we need to multiply both sides by $-1$ to get:

\begin{center}
    $v' - \frac{4}{x} v = -x^3$ \\
\end{center}

Next, we need to multiply all variables by $\mu(x)$

\begin{center}
    $\mu(x)v' - \mu(x)\frac{4}{x} v = -\mu(x)x^3$ \\
\end{center}

next, we need to find $\mu(x)$

$\mu(x) = e^{\int -\frac{4}{x} = -4\int\frac{1}{x} = -4 ln |x| \rightarrow ln|x^{-4}} = x^{-4}$

We can plug this in to get

\begin{center}
    $x^{-4}v' - x^{-4}\frac{4}{x} v = -x^{-4}x^3$ = $-x^4v' - 4x^{-5}v = -x^{-1}$ \\
\end{center}

You can see that this resembles the product rule, hence we can simplify this equation to:

$-x^{-4}x^3$ = $-x^4v' - 4x^{-5}v = (x^{-4}v)'$

We can then integrate:

\begin{center}
    $\int (x^{-4}v)' = \int -x^{-1}dx$ \\
    $x^{-4}v = -ln|x| + c$ \\
    $v(x) = cx^4 - x^4 ln x$, or $v = cx^4 - x^4 ln x$
\end{center}

We can then solve for c by using the condition $y(2) = -1$. We can also plug $y^-1 back in$

\begin{center}
    $y^{-1} = cx^4 - x^4 ln x$ \\
    $-1^{-1} = c2^4 - 2^4 ln 2$ \\
    $-1 + 16 ln 2 = 16c - 16 ln 2 + 16 ln 2$ \\
    $\frac{-1 + 16 ln 2 }{16}= \frac{16c}{16}$
    $c = \frac{-1 + 16ln 2}{16}$

\end{center}

We can solve for y, bla bla bla. 

\section{Substitutions}
\subsection{Introduction}

In addition to Bernoulli's Differential equation, there's also another one known as homogenous differential equations.
\subsubsection{Homogenous equations}

Homogenous equations are equations that take the form:

\begin{center}
    $y' = F(\frac{y}{x})$
\end{center}

If this is the case, we can use the following substitutions:

\begin{center}
    $v = \frac{y}{x}$, which can be written as $y = xv$\\
    $y' = v + xv'$
\end{center}

using this will give us a separable differential equation, which we can easily solve. 

\paragraph{Example}

Solve the following IVP and find the interval of validity for the solution:

\begin{center}
    $xyy' + 4x^2 + y^2 = 0$, where $y(2) = -7$ and $x > 0$
\end{center}

We can move $4x^2 + y^2$ to the other side of the equation, giving us:

\begin{center}
    $xyy' = -4x^2 - y^2$
\end{center}

we can then divide both sides by two:

\begin{center}
    $\frac{y}{x} y' = -4 - \frac{y^2}{x^2} = -4 - (\frac{y}{x})^2$
\end{center}
Hence, we proved our equation here is homogenous.  Next, we can apply the substitution here.

\begin{center}
    $v(v + xv') = -4 - v^2$
\end{center}

We can then separate variables, then use integration to solve our differential equation.

\begin{center}
    $\int \frac{v}{2v^2} dv = \int -\frac{1}{x}dx$
\end{center}

\subsubsection{Other equation}

If you have a differential equation with the form: $y' = G(ax + by)$, then you can apply the following substitution:

\begin{center}
    $v = ax + by$ \\
    $v' = a + by'$
\end{center}
\section{Intervals of Validity}
\subsection{Theorem one}
When it comes to intervals of validity, there's an important idea:

\begin{center}
    $y' + p(t)y = g(t)$, where $y(t_0) = y_0$
\end{center}

If $p(t)$ and $g(t)$ are continuous functions on an open interval $\alpha < t < \beta$, containing an interval $t_0$, then there's a unique solution to the IVP on the interval.

\subsubsection{Example}
Let's try to determine the intervals of validity for the following initial value problem.

\begin{center}
    $(t^2 - 9)y' + 2y = ln|20 - 4t|$, where $y(4) = -3$
\end{center}

Let's start off by dividing both sides by $(t^2 - 9)$, giving us:

\begin{center}
    $y' + \frac{2}{t^2 - 9} y = \frac{ln|20 - 4t|}{t^2 - 9}$
\end{center}

In order to find non continuous intervals, we need to find points that aren't defined. In this example, dividing by zero, or if you get $ln|0|$.

\begin{center}
    $20 - 4t = 0 \rightarrow t = 5$ \\
    $t^2 - 9 = 0 \rightarrow (t - 3)(t + 3) = 0 \rightarrow t = \pm 3$
\end{center}

hence, our Intervals of Validity include:

\begin{center}
    $-\infty < t < -3, -3 < t < 3, 3 < t <  5, 5 < t < \infty$
\end{center}

\subsection{Theorem 2}

Let's consider the following IVP:

\begin{center}
    $y' = f(t,y)$, where $y(t_0) = y_0$
\end{center}

if $f(t,y)$ and $\frac{\partial f}{\partial y}$ are continious functions in some rectangle $\alpha < t < \beta$, and $\gamma  < y < \delta$, containing the point $(t_0,y_0)$, then there's a unique solution to the IVP in some interval $t_0 - h < t < t_0 + h$

\subsubsection{Example}

Determine the interval of validity for the initial value problem below and give its dependence on the value of $y_0$

\begin{center}
    $y' = y^2$, where $y(0) = y_0$
\end{center}

Here, we can define the variables:

\begin{itemize}
    \item $f(t,y) = y^2$
    \item $t_0 = 0$
    \item $y_0$ is the initial value.
\end{itemize}

We first need to check for continuity for $f(t,y)$ and $\frac{\partial f}{\partial y}$
\begin{itemize}
    \item $f(t,y) = y^2$ is continuous
    \item $\frac{\partial f}{\partial y} = 2y$, which is also continuous.
\end{itemize}

Our equation also contains the point $(t_0,y_0)$, so we're good. we can now 

now, we can work on this:

\begin{center}
    $\frac{dy}{dt} = y^2$ \\
    $dy = y^2 dt$\\
    $\frac{1}{y^2} dy = dt$ \\
    $\int \frac{1}{y^2} dy = \int dt$ \\
    $-\frac{1}{y} = t + c$
\end{center}

We're given $y(0) = y_0$, so we can figure out c with it.

\begin{center}
    $- \frac{1}{y_0} = c$
\end{center}

therefore, our solution is $-\frac{1}{y} = t - \frac{1}{y_0}$.

We can rearrange the variables to get $y = \frac{y_0}{1 - y_0 t}$

We can then find out that y is undefined when $t = \frac{1}{y_0}$.

\part{Second Order Differential Equations}

\section{Basic Concepts}

\subsection{Second order linear differential equation}
The most general linear order differential equation follows the form:

\begin{center}
    $P(t)y'' + q(t)y' + r(t)y = g(t)$
\end{center}

If we were to assume constant coefficient, we can use the following:

\begin{center}
    $ay'' + by' + cy = g(t)$
\end{center}

\subsection{Homogenous vs. nonhomogenous differential equations}

Given the form, $ay'' + by' + cy = g(t)$:
\begin{itemize}
    \item we say a differential equation is homogenous if $g(t) = 0$
    \item Otherwise, if $g(t) \neq 0$, it's nonhomogenous.
\end{itemize}

\subsubsection{Example}

Determine some solutions to $y'' - 9y = 0$

We can bring 9y to the other side:

\begin{center}
    $y'' = 9y$
\end{center}

\begin{itemize}
    \item This tells us $y''$ needs to be the original function times 9.
    \item Right off the bat, we can think of using e.
    \item more specifically, $y = e^{\pm 3t}$
\end{itemize}

We can test this:

\begin{center}
    $y = e^{3t}$ \\
    $y' = 3e^{3t}$ \\
    $y'' = 9e^{3t}$ \\
\end{center}

since $y = e^{3t}$, we can say $y'' = 9y$ 

\subsection{Principle of Superposition}

If $y_1 (t)$ and $y_2 (t)$ are solutions to a linear, homogenous differential equation, we can also say $y(t) = c_1 y_1 (t) + c_2 y_2 (t)$ is a solution to our differential equation. \textbf{This is true regardless of order}, applies to any homogenous linear differential equation.

Another important question to ask is, how do we find the constants, $c_1$ and $c_2$? We simply just specify values of the solutions at two distinct points.

\subsubsection{Example}

Solve the following IVP

\begin{center}
    $9'' - 9y = 0$, where $y(0) = 2$ and $y'(0) = -1$
\end{center}

\begin{itemize}
    \item We solved y in the first part, it's $e^{\pm 3t}$
    \item Let's plug it in: $y(t) = c_1 e^{3t} + c_2 e^{-3t}$
    \item next, we can figure out the derivative of $y(t)$
    \begin{itemize}
        \item $y'(t) = 3c_1 e^{3t} - 3c_2 e^{-3t}$
    \end{itemize}
    \item We can plug in our initial conditions:
    \begin{itemize}
        \item $2 = c_1 + c_2$, which is $y(0)$
        \item $-1 = 3c_1 - 3c_2$, which is $y'(0)$
    \end{itemize}
    \item We can use gaussian elimination to find the t$c_1 = \frac{5}{6}$ and $c_2 = \frac{7}{6}$
    \item Plugging everything in, we can find that our differential equation is $y(t) = \frac{5}{6} e^{3t} + \frac{7}{6} e^{-3t}$
\end{itemize}

\subsection{Characteristic equations}

\begin{itemize}
    \item We can notice that the equations above resembled $y(t) = e^{rt}$.
    \item How can we solve without guessing?
\end{itemize}

\subsubsection{Example}

Let's do the $y'' - 9y = 0$

Let's first calcule y, which we'll say is $y(t) = e^rt$

\begin{itemize}
    \item $y(t) = e^{rt}$
    \item $y'(t) = re^{rt}$
    \item $y''(t) = r^2e^{rt}$
\end{itemize}

We can plug these into our equation:

\begin{center}
    $r^2 e^{rt} - 9e^{rt} = 0$
\end{center}

from here, we can bring the $e^{rt}$ out the equation:

\begin{center}
    $e^{rt}(r^2 - 9) = 0$
\end{center}

\begin{itemize}
    \item $r^2 - 9$ is known as a \textbf{characteristic equation}.
    \item Based on this, we can deduce that, in order for the equation to be equal to zero, $r^2 - 9 = 0$. We can use easy algebra to find that $r = 3$, and $r = -3$.
\end{itemize}

\subsection{Real \& distinct roots}

Let's think about our characteristic equation:

\begin{center}
    $ar^2 + br + c = 0$
\end{center}

and the two solutions: $y_1(t) = e^{r_1 t}$ and $y_2(t) = e^{r_2 t}$
\begin{itemize}
    \item We say that the roots, $r_1$ and $r_2$ are real and distinct, if the two are real numbers and $r_1 \neq r_2$.
    \item If this is true, then the two solutions are \textbf{nice enough} to form the solution: $y(t) = c_1 e^{r_1 t} + c_2 e^{r_2 t}$
\end{itemize}

\subsubsection{Example}

Solve the following IVP:

\begin{center}
    $y'' + 11y' + 24y = 0$, given $y(0) = 0$, and $y'(0) = -7$.
\end{center}

Let's use $y = e^{rt}$

\begin{center}
    $y = e^{rt}$ \\
    $y' = re^{rt}$ \\
    $y'' = r^2 e^{rt}$
\end{center}

Let's plugs these in our equation:

\begin{center}
    $r^2 e^{rt} + 11re^{rt} + 24e^{rt} = 0$ \\
    $e^{rt}(r^2 + 11r + 24) = 0$ \\
    $r^2 + 11r + 24 = 0$ \\
    $(r + 8)(r + 3) = 0$ \\
    therefore, $r = -8$ and $r = -3$, which are real numbers and distinct.
\end{center}

we can now come up with the following equations:

\begin{center}
    $y(t) = c_1 e^{-3t} + c_2 e^{-8t}$ \\
    $y'(t) = -3c_1^{-3t} - 8c_2 e^{-8t}$
\end{center}

We can use gaussian elimination to solve for the constants. Our final answer is $y(t) = -\frac{7}{5}e^{-3t} + \frac{7}{5} e^{-8t}$

\subsection{Complex Roots}
\subsubsection{Introduction}

Let's again, look at solutions to the differential equations following the form:

\begin{center}
    $ay'' + by' + cy = 0$ \\
    $ar^2 + br + c = 0$
\end{center}

We call roots \textbf{complex roots}, if they have the form $r_{1,2} = \lambda \pm \mu i$.

Assuming we're using the two differential equations, $y(t) = e^{rt}$, we say that our solutions to the differential equation is: 

\begin{center}
    $y_1 (t) = e^{(\lambda + \mu i) t}$ \\
    $y_1 (t) = e^{(\lambda - \mu i) t}$
\end{center}

Having complex roots is great, but it's ideal for our solution to only have real numbers. we use \textbf{Euler's formula} for this:

\begin{center}
    $e^{i\theta} = cos\theta + isin\theta$
\end{center}

We can use another variant:

\begin{center}
    $e^{-i\theta} = cos(-\theta) + isin(-\theta) = cos\theta - isin\theta$
\end{center}

We can now split our exponent into two groups:

\begin{center}
    $y_1(t) = e^{\lambda t + \mu i t} = e^{\lambda t} e^{\mu i t }$
    $y_2(t) = e^{\lambda t - \mu i t} = e^{\lambda t} e^{-\mu i t }$
\end{center}

We have our equation from above, which we can use:

\begin{center}
    $y_1(t) = e^{\lambda t} e^{i\mu t} = e^{\lambda t}(cos(\mu t) + isin(\mu t))$
    $y_2(t) = e^{\lambda t} e^{-i\mu t} = e^{\lambda t}(cos(\mu t) - isin(\mu t))$
\end{center}

We know that, given two solutions are "nice enough", any solution can be written as:

\begin{center}
$y(t) = c_1 y_1 (t) + c_2 y_2(t)$
\end{center}

We can ignore constants for now, so 

\begin{center}
    $y(t) = y_1(t) + y_2(t)$ \\
    $y_1(t) + y_2(t) = e^{\lambda t}cos(\mu t) + e^{\lambda t} i sin(\mu t) + e^{\lambda t} cos(\mu t) - e^{\lambda t}isin(\mu t)$ \\
    $y_1(t) + y_2(t) = e^{\lambda t}cos(\mu t) + e^{\lambda t} cos(\mu t)$  \\
    $y_1(t) + y_2(t) = 2(e^{\lambda t} cos(\mu t))$ 
\end{center}

Now, we can assume the constants, $c_1$ and $c_2$, are $\frac{1}{2}$, which means

\begin{center}
    $\frac{1}{2}y_1(t) + \frac{1}{2}y_2(t) = (e^{\lambda t} cos(\mu t))$ 
\end{center}

\textbf{main idea}

We can say, if the roots of the characteristic equation happens to be $r_{1,2} = \lambda \pm \mu i$, the solution to the differential equation is:

\begin{center}
    $y(t) = c_1 e^{\lambda t}cos(\mu t) + c_2 e^{\lambda t}sin(\mu t)$
\end{center}

\subsubsection{Example}

Solve the following IVP:

\begin{center}
    $y'' - 4y' + 9y = 0$, where $y(0) = 0$ and $y'(0) = -8$
\end{center}

our characteristic equation is $r^2 - 4r + 9 = 0$. There's no easy way to solve this, so we need to use the quadratic formula.

\begin{center}
    $r = \frac{-b \pm \sqrt{b^2 -4ac}}{2a}$
\end{center}

Where $a = 1$, $b = -4$ and $c = 9$, which we can plug in:

\begin{center}
    $r = \frac{4 \pm \sqrt{-4^2 -4(1)(9)}}{2(1)}$, which simplifies to: \\
    $r = 2 \pm \sqrt{5}i$
\end{center}

This resembles our complex root equation, $\lambda \pm \sqrt{\mu}i$, where $\lambda = 2$ and $\mu = 5$. We can plug these variables into our equation

\begin{center}
    $y(t) = c_1 e^{2 t}cos(\sqrt{5} t) + c_2 e^{2 t}sin(\sqrt{5} t)$
\end{center}

we can now think about our first initial condition, $y(0) = 0$


\begin{center}
    $0 = c_1 e^{2(0)}cos(\sqrt{5} (0)) + c_2 e^{2 (0)}sin(\sqrt{5} (0))$ \\
    $0 = c_1 \times 1$ \\
    $c_1 = 0$
\end{center}

this cancels out the first part, giving us:

\begin{center}
    $y(t) = c_2 e^{2 t}sin(\sqrt{5} t)$
\end{center}

now, we need to find the derivative, which is done using the product rule:

\begin{center}
    $y'(t) = c_2 e^{2t} \sqrt{5}cos(\sqrt{5} t) + 1c_2 e^{2t} sin(\sqrt{5} t)$
\end{center}

Here, we can use the initial condition, $y'(0) = -8$

\begin{center}
    $-8 = c_2 e^{2(0)} \sqrt{5}cos(\sqrt{5} (0)) + 1c_2 e^{2(0)} sin(\sqrt{5} (0))$ \\
    $-8 = c_2\sqrt{5}$ \\
    $c_2 = -\frac{8}{\sqrt{5}}$
\end{center}

Therefore, our solution is:

\begin{center}
    $y(t) = -\frac{8}{\sqrt{5}} e^{2 t}sin(\sqrt{5} t)$
\end{center}

\subsection{Repeated Roots}
\subsubsection{Introduction}
Here, if the roots are the same, or $r_1 = r_2 = r$, we can say the general solution will be $y(t) = c_1 e^{rt} + c_2 te^{rt}$

\subsubsection{Example}

Solve the following IVP:

\begin{center}
    $y'' - 4y' + 4y = 0$, where $y(0) = 12$ and $y'(0) = -3$
\end{center}

We can get the characteristic equation

\begin{center}
    $r^2 - 4r + 4 = 0$ \\
    $(r - 2)(r - 2) = 0$ \\
    $r_1 = r_2 = 2$
\end{center}

We can see that this falls under the Repeated root, therefore, we can plug r into the general solution:

\begin{center}
    $y(t) = c_1 e^{2t} + c_2 te^{2t}$
\end{center}

Now, we can solve for the constants, first using $y(0) = 12$:

\begin{center}
    $12 = c_1 e^{2(0)} + c_2 (0)e^{2(0)}$ \\
    $12 = c_1$
\end{center}

We can now find the derivative of $y(t)$, and finde $c_2$ with $y'(0) = -3$

\begin{center}
    $y(t) = 12 e^{2t} + c_2 te^{2t}$ \\
    $y'(t) = 24 e^{2t} + c_2 e^{2t} + 2c_2te^{2t}$ (used product rule) \\
    $-3 = 24 e^{2(0)} + c_2 e^{2(0)} + 2c_2(0)e^{2(0)}$  \\
    $-3 = 24 + c_2$ \\
    $-27 = c_2$
\end{center}


Therefore, our solution is $y(t) = 12 e^{2t} - 27 te^{2t}$

\end{document}

